\section{Alarm Mechanism and Gap Decomposition}
\label{sec:gap_decomposition}

\subsection{The alarm}

\begin{definition}[B-to-C Evidence Alarm]
\label{def:alarm}
The B-to-C evidence alarm fires when the lower-bound B-to-C
ratio drops below a threshold:
\[
  \mathrm{ALARM}(T) \;=\; \betaL(P, T) \;<\; \betaalarm,
\]
where $\betaL(P, T)$ is the global conservative ratio of the
companion paper~\citep[Revealed~Sacrifice]{lasser2026paper8a} on the full
capability poset $P$ (numerator $\volRlower(U, T)$ extended by
zero outside the observed subset $U$, denominator $\volL(P)$),
and $\betaalarm \in (0, 1)$ is configurable.  The
gap-decomposition identity below is stated against this same
denominator $\volL(P)$; the alarm and decomposition use a
common denominator throughout.  The alarm is an
\emph{evidence alarm}, not a proxy-failure alarm directly: when
$\betaL$ is low, the sacrifice channel shows little evidence of
$\volR$-content proportional to $\volL$, but this does not by
itself imply that true exercise is low.  The interpretation of a
firing alarm requires additional analysis (via the gap
decomposition below) to determine whether low evidence reflects
genuine low exercise, restricted capabilities, dormant wants,
or residual-class capabilities beyond the channel's reach.
\end{definition}

\begin{remark}[Evidence alarm versus proxy-failure diagnosis]
\label{rem:alarm_interpretation}
Let $\betatrue := \volR/\volL$ be the true (latent) exercised
fraction, of which $\betaL = \volRlower/\volL$ is the
sacrifice-channel lower bound.  Because $\betaL \leq \betatrue$
by construction, the alarm distinguishes two regimes by which
$\delta$ term dominates the gap $1 - \betaL$:
\begin{itemize}
\item \textbf{Alarm fires and $\delta_{\mathrm{covered}}$
  dominates the gap:} on the covered cell, observed trade
  evidence is small relative to that cell's $\volL$-weight
  ($\delta_{\mathrm{covered}} = [\volL(C) - \volRlower(C)] /
  \volL(P)$ is large, indicating
  $\volRlower(C) \ll \volL(C)$).  This is the
  proxy-failure-\emph{candidate} signature: capabilities the
  sacrifice channel can reach are present in $\volL$ but
  registering little exercise evidence.  Reaching a
  proxy-failure conclusion from this signature additionally
  requires ruling out the lower-bound's slack as the
  explanation: low $\volRlower(C)$ relative to $\volL(C)$ can
  reflect (i)~genuine low exercise (the proxy-failure
  reading), (ii)~Part-A supremum or Part-B max-attribution
  loose against the true $\volR(C)$ (the lower-bound is
  conservative by construction, and its conservatism varies
  with bundle structure and the $\alpha_{n,c}$ defaults),
  (iii)~admissibility filtering rejecting events that would
  otherwise contribute to $\volRlower(C)$, or (iv)~sparse
  observation in $W$.  The covered-deficit signature is
  necessary but not sufficient for a proxy-failure
  determination; institutional review of the explanation
  is required before escalation.
\item \textbf{Alarm fires and $\delta_{\mathrm{residual}}$ or
  $\delta_{\mathrm{restricted}}$ dominates:} the low $\betaL$
  reflects structural invisibility or policy restriction, not
  proxy failure.  The framework should not escalate in these
  cases; the alarm is firing on what the sacrifice channel cannot
  reach, not on degraded exercise on what it can reach.
\end{itemize}
As a diagnostic for genuine
$\volR$-degradation, the alarm has both
false-positive and false-negative regimes:
\begin{itemize}
\item \emph{False positives} (alarm fires, no proxy failure)
  arise when low $\betaL$ reflects channel-blindness (high
  $\delta_{\mathrm{residual}}$ or $\delta_{\mathrm{restricted}}$),
  bound looseness (Part-A's supremum or Part-B's
  max-attribution falling well below the truth on covered
  capabilities), admissibility-filter exclusions, or sparse
  observation rather than degraded exercise.  The gap
  decomposition's cell-attribution is the diagnostic that
  distinguishes these from genuine proxy failure.
\item \emph{False negatives} (alarm does not fire, but
  proxy-failure pathology is present) arise when wireheading
  or other pathologies produce high $\volRlower$ on
  loop-driver categories without depressing $\betaL$ below the
  threshold; the wireheading-consistent concentration signal
  of Section~\ref{sec:wireheading} is the diagnostic
  complementing $\betaL$ for these cases.  Pathologies that
  neither lower $\betaL$ nor concentrate observed flow are
  outside the joint signal's reach.
\end{itemize}
Both regimes are structural costs of privacy-respecting
observation; the multi-diagnostic combination
(Proposition~\ref{prop:combined_diagnostic}) is what reduces
their individual blind spots.
\end{remark}

\subsection{Diagnostic decomposition}

\begin{definition}[Gap Decomposition Under Revealed Sacrifice]
\label{def:gap_decomp}
Under the Part~B hypotheses of the aggregate lower bound
theorem of the companion paper~\citep[Revealed~Sacrifice]{lasser2026paper8a}
(S0--S5), the per-capability aggregate $\volRlower(c)$ is
well-defined, and subset aggregates $\volRlower(A)$ are given
by the subset-extension of $\volRlower$~\citep[Revealed~Sacrifice]{lasser2026paper8a}.  Under those
hypotheses, the B-to-C gap ($1 - \betaL$) decomposes into
\emph{disjoint} sources whose count depends on the closure
strategy (Remark~\ref{rem:cell_cooperative_closure}): five
under the default literal-partition strategy~(c), four under
cooperative-closure strategies~(a)/(b).  Each individual
capability in $P$ is assigned to exactly one cell by priority
ordering (restricted $>$ covered $>$ dormant $>$ residual);
under~(c), boundary-crossing cooperative nodes form a fifth
bucket~$B$.  Each cell's $\delta$ term is its uncovered $\volL$
share:
\[
  1 - \betaL \;=\; \delta_{\mathrm{restricted}}
  + \delta_{\mathrm{covered}} + \delta_{\mathrm{dormant}}
  + \delta_{\mathrm{residual}}
  + \delta_{\mathrm{boundary}},
\]
where $\delta_{\mathrm{boundary}}$ is omitted under
strategies~(a)/(b) (which absorb cross-cell cooperatives via
priority-extension into a participant's cell rather than
routing them to~$B$).  The capability partition is constructed
in two stages.
\medskip\noindent
\textbf{Stage 1 (raw seed cells, by priority).}
Classify individual capabilities by the priority rule:
\begin{itemize}
\item \textbf{Seed restricted cell:} $R_0 := \Xi$, where $\Xi$
  is the set of capabilities currently restricted under the
  preemptive-restriction criterion
  of~\citep[Proposition~1]{lasser2026horizon}.  Classified first
  regardless of sacrifice evidence.
\item \textbf{Seed covered cell:} $C_0 := \CovS \setminus \Xi$,
  where $\CovS$ is the set of capabilities appearing in at
  least one acquired bundle $Y_n$ of a theorem-admissible
  sacrifice event in $[T - H, T]$.  A theorem-admissible event
  for Part~A satisfies S0, S1, S2, and S4(a); for Part~B it
  additionally satisfies S3, S5, and within-bundle
  poset-independence (the event-local bundle
  decomposition~\citep[Revealed~Sacrifice]{lasser2026paper8a}).
\item \textbf{Seed dormant cell:} $D_0 := \DorS$, where
  $\DorS = \{d \in P : d \notin \Xi,\; d \notin \ResS,\;
  d \notin \CovS\}$.  Tradable-in-principle capabilities with
  zero want-pursuit sacrifice evidence---dormant \emph{wants}.
\item \textbf{Seed residual cell:} $S_0 := \ResS \setminus \Xi$,
  where $\ResS$ is the structural residual class
  (Definition~\ref{def:residual_class}).  Residual capabilities
  outside $\Xi$; capabilities in $\Xi \cap \ResS$ are already
  assigned to $R_0$ by priority.
\end{itemize}
The four seed cells $\{R_0, C_0, D_0, S_0\}$ are pairwise
disjoint and their union is~$P$: $R_0$ and $S_0$ are disjoint by
construction; $C_0$ and $D_0$ are priority-filtered to exclude
$\Xi$; $C_0 \cap \ResS = \emptyset$ because residual
capabilities admit no sacrifice-observable path in principle
(Definition~\ref{def:residual_class}) and therefore cannot
appear in any admissible bundle; remaining pairs are disjoint
by the bullet-stated exclusions.

\medskip\noindent
\textbf{Stage 2 (closure-adjusted cells).}
The seed cells may contain cooperative or subsumption edges
crossing cell boundaries; to obtain a closed additive identity
we apply a closure policy.  The default is the literal-partition
variant (strategy~(c) of
Remark~\ref{rem:cell_cooperative_closure}), which preserves
seed-cell attribution by routing every boundary-crossing
cooperative node into a separate boundary-residual bucket~$B$
rather than absorbing it into a participant's cell; this yields
the five-cell partition $\{R, C, D, S, B\}$ used by the worked
example, with the additivity identity following from the
strategy-(c) atomic-weight accounting convention
$\volL(X) := \sum_{c \in X} w(c)$ summed over the disjoint
cells (no M4 invocation; see
Remark~\ref{rem:cell_cooperative_closure} for the precise
semantics).
Under cooperative-closure semantics
(strategies~(a)/(b)), an alternative four-cell partition
$\{R, C, D, S\}$ is available at the cost of priority-extension
absorbing seed-dormant and seed-residual atoms into the covered
cell when they are poset-connected to a non-residual node.  In
either regime the cells are pairwise
disjoint, union to~$P$, and agree with the seeds on all
individual-capability classifications except where closure
redistributes poset-connected components
(Remark~\ref{rem:residual_cell_irreducibility} discusses when
the absorption occurs under (a)/(b) and the condition under
which the residual cell preserves its irreducibility).  The
four-cell variant additionally satisfies:
\begin{enumerate}
\item[(CC1)] \emph{Cooperative closure.}  For every cooperative
  capability $z \in P$, all participants of $z$ lie in the same
  cell.
\item[(CC2)] \emph{Subsumption closure.}  No subsumption edge
  $c \succeq c'$ crosses a cell boundary: if $c, c'$ are
  connected by subsumption in $P$, they lie in the same cell.
\end{enumerate}
Under the \emph{cooperative-closure semantics} for cell
$\volL$ values---where each cell $X$ is read as $X^{*} := X
\cup \{z : z \text{ cooperative with all participants in } X\}$
and $\volL(X)$ is summed over $X^{*}$---set-disjointness alone
does not suffice for $\volL$-additivity across cells.  By the
category-partition poset-independence
condition of the companion paper~\citep[Revealed~Sacrifice]{lasser2026paper8a} and
the M4 axiom of~\citep[Proposition~1]{lasser2026poset},
additivity requires full poset-independence---\emph{neither}
cross-cell cooperative composition \emph{nor} cross-cell
subsumption.  A cooperative
whose participants straddle two cells contributes to $\volL(P)$
but to neither cell's closure individually (CC1 failure); a
subsumption edge crossing the boundary couples $\volL$-mass
across cells in ways M4's independent-component additivity
cannot decompose (CC2 failure).  Together (CC1) and (CC2) give
the poset-independence that licenses $\volL(P) = \sum_X
\volL(X)$ for $X \in \{R, C, D, S\}$ under the
cooperative-closure semantics.  Strategy~(c) below uses a
\emph{literal-partition semantics} instead, where each
capability (cooperatives included) is assigned to exactly one
cell by stipulation and cell-level $\volL$ values are read as
atomic-weight accounting aggregates
$\volL(X) := \sum_{c \in X} w(c)$ rather than as the M1--M6
poset measure on poset-independent components; under that
weaker semantics, the identity follows from direct rearrangement
of the atomic sum across set-disjoint cells without the
(CC1)/(CC2) hypotheses, but the cells are no
longer M4-additive in the cooperative-closure sense.  See
Remark~\ref{rem:cell_cooperative_closure} for both regimes.

The $\delta$ terms are:
\begin{align*}
  \delta_{\mathrm{restricted}} &= [\volL(R) -
    \volRlower(R)] \;/\; \volL(P), \\
  \delta_{\mathrm{covered}} &= [\volL(C) -
    \volRlower(C)] \;/\; \volL(P), \\
  \delta_{\mathrm{dormant}} &= \volL(D) \;/\; \volL(P), \\
  \delta_{\mathrm{residual}} &= \volL(S) \;/\; \volL(P), \\
  \delta_{\mathrm{boundary}} &= [\volL(B) -
    \volRlower(B)] \;/\; \volL(P)
    \quad \text{(default strategy~(c) only;}
    \\ &\hphantom{={}} \text{omit under~(a)/(b))}.
\end{align*}
Each $\delta \geq 0$: $\volRlower(X) \leq \volR(X) \leq
\volL(X)$ holds for any cell~$X$ (the first inequality by
construction of $\volRlower$ as a lower bound on $\volR$
\citep[Revealed~Sacrifice]{lasser2026paper8a}, the second because $\volR$ is
$\volL$ restricted to the exercised sub-poset and so cannot
exceed $\volL$ on any subset), so the
$\delta_{\mathrm{restricted}}$ and $\delta_{\mathrm{covered}}$
numerators are non-negative.  The $\delta_{\mathrm{dormant}}$
and $\delta_{\mathrm{residual}}$ terms are non-negative by
M5 ($\volL \geq 0$).
Two additivity stories underlie the identity, depending on
closure strategy: (a)/(b) use M4 on poset-boundary-closed cells
to obtain $\volL$ additivity, while (c) uses literal weight
summation on a five-bucket set partition.  These are not
incompatible: (c) is a strictly weaker license that delivers
the identity by direct summation without invoking M4 on
cells, while (a)/(b) deliver the same identity via M4 on the
four poset-boundary-closed cells.  Under strategies (a) or (b), the
identity $1 - \betaL = \sum \delta$ follows from two additivity
steps.  First, $\volL(P) = \volL(R) + \volL(C) + \volL(D) +
\volL(S)$ by poset-boundary closure plus set-disjoint union via
M4~\citep[Proposition~1]{lasser2026poset}: CC1 prevents any
cooperative from spanning two cells and CC2 prevents any
subsumption from crossing, so the cells form an M4-additive
partition for $\volL$.  Second,
$\volRlower(P) = \volRlower(R) + \volRlower(C) + 0 + 0$ by
the Part-B subset-extension of $\volRlower$, which under
Part-B's per-capability max-attribution
\citep[Revealed~Sacrifice]{lasser2026paper8a}
is defined by $\volRlower(A) = \sum_{c \in A} \volRlower(c)$,
making set-disjointness alone sufficient for additivity across
cells; $D$ and $S$ contribute zero by definition (no admissible
sacrifice evidence).  CC1 and CC2 are
\emph{not} required for this $\volRlower$ step---they are
required only for the $\volL$ step above.  Under the literal-partition
variant (strategy~(c)), the $\volL$ identity expands to a
five-term sum $\volL(P) = \volL(R) + \volL(C) + \volL(D) +
\volL(S) + \volL(B)$ by literal M5 weight summation across the
five-cell partition; the $\volRlower$ aggregate similarly
sums literally over the partitioned subset.  $\volRlower(B)$
is computed by the same Part-B max-attribution rule as any
other cell---it is \emph{not assumed zero}---and equals zero
only when no admissible bundle covers a boundary cooperative
in $B$.  This is the typical case (cross-cell cooperatives
whose participants span at least one dormant or residual cell
inherit zero $\volRlower$ contribution from those cells'
absence of admissible bundle coverage), but it is a
typical-case observation, not a definitional one; a
cross-cell cooperative whose participants happen to be covered
by an admissible bundle that contains the cooperative as a
whole would contribute non-zero $\volRlower(B)$, and the
$\delta_{\mathrm{boundary}}$ formula
$[\volL(B) - \volRlower(B)]/\volL(P)$ accommodates this
case.  Combining either route:
$1 - \betaL = 1 - \volRlower(P)/\volL(P) = [\volL(P) -
\volRlower(P)] / \volL(P) = \sum_X [\volL(X) -
\volRlower(X)] / \volL(P)$, which expands to the four-term
identity under (a)/(b) (with $\volRlower(D) = \volRlower(S) =
0$ by the dormant/residual definitions, $B = \emptyset$) or the
five-term identity under (c) (with $\volRlower(D) =
\volRlower(S) = 0$ and $\volRlower(B)$ given by the
max-attribution rule, typically zero).
\end{definition}

\begin{remark}[When the residual cell preserves its
irreducibility under closure]
\label{rem:residual_cell_irreducibility}
Under the default literal-partition strategy~(c) the residual
cell~$S$ always equals the seed $S_0 = \ResS \setminus \Xi$
because (c) does not redistribute seed-classified atoms
across cells (boundary-crossing cooperative nodes go to~$B$,
not to a participant's cell).  The remainder of this remark
addresses the cooperative-closure strategies~(a)/(b), where
seed-cell preservation is conditional.  Under~(a)/(b),
the closure-adjusted residual cell~$S$ equals the seed
$S_0$ when $\ResS$ is itself
\emph{poset-isolated}: no cooperative or subsumption edge
connects a residual capability to a non-residual one.  Under
priority-extension, an edge from $c \in
\ResS \setminus \Xi$ to a non-residual node $c'$ in $C_0$,
$D_0$, or $\Xi$ would absorb the entire connected component
into the higher-priority cell of $c'$ (residual having lowest
priority); the absorbed mass then registers as
$\delta_{\mathrm{covered}}$, $\delta_{\mathrm{dormant}}$, or
$\delta_{\mathrm{restricted}}$ rather than as
$\delta_{\mathrm{residual}}$.  The
``residual is structurally irreducible'' interpretation of
Definition~\ref{def:residual_class} therefore holds for the
$S$-cell of the decomposition only when~$\ResS$ is
poset-isolated; when it is not, residual mass migrates to
neighbouring cells under priority-extension.  Operationally,
$\ResS$ being poset-isolated is the typical case because the
residual class characterizes purely-private exercise
(Definition~\ref{def:residual_class}), which by construction
does not participate in cooperative or subsumption structure
with publicly-observable capabilities; verifying isolation is
nonetheless an explicit residual-class hygiene check.
Strategies (b) and (c) of
Remark~\ref{rem:cell_cooperative_closure} preserve residual
mass by other means: (b) sub-partitions cells by connected
component, so residual components are accounted for separately;
(c) sends boundary-crossing mass to a fifth bucket rather than
absorbing it into a neighbouring cell.
\end{remark}

\begin{remark}[Achieving poset-boundary closure of the cells]
\label{rem:cell_cooperative_closure}
Three operational strategies admit the poset-boundary closure
hypothesis without changing the framework's substantive content.
Each strategy handles both the cooperative (CC1) and subsumption
(CC2) conditions jointly, since both are instances of the
``edge connects different cells'' pattern; concretely, treat the
poset's combined cooperative/subsumption graph (each cooperative
contributes hyperedges to its participants; each subsumption
relation contributes a directed edge) as the object whose
boundary-crossing instances must be eliminated.
\begin{enumerate}
\item[(a)] \emph{Priority-extension.}  For each node
  $n \in P$ not yet in a settled cell, and for each subsumption
  edge $c \succeq c'$ or cooperative $z$ with participant set
  $\Pi(z)$ touching $n$: take the transitive closure under
  cooperative participation and subsumption edges, collect the
  resulting component, and place the whole component in the
  cell of its highest-priority member
  (restricted $>$ covered $>$ dormant $>$ residual).  After one
  sweep, the cells are both cooperative-closed and
  subsumption-closed (hence poset-boundary closed), and the
  original individual-capability classification is preserved
  except where a priority-dominated node was absorbed into a
  higher-priority cell.  \emph{Caveat: dormant absorption.}
  Priority-extension can absorb dormant or residual atomic
  capabilities into the covered cell when they are
  poset-connected to a covered atom.  Their $\volL$ weight
  then contributes to $\volL(C)$ but not to $\volRlower(C)$
  (since they had no admissible sacrifice evidence by hypothesis),
  inflating $\delta_{\mathrm{covered}}$ relative to a literal
  classification.  This means \emph{$\delta_{\mathrm{covered}}$
  under priority-extension does not exclusively measure
  proxy-failure on observed capabilities}; some of its mass
  may be absorbed dormancy or absorbed residual.  Strategies
  (b) and (c) do not have this absorption effect, at the cost
  of weaker M4 license.
\item[(b)] \emph{Poset-closed sub-partition.}  Refine the
  partition so each cell is sub-partitioned into its
  \emph{poset-connected components} under both cooperative
  composition and subsumption \emph{globally across the entire
  $P$}, not within the original cell only: the decomposition
  operates over the global connected-component partition of
  $P$, then assigns each component to a cell by its
  highest-priority member (as in strategy~(a)).  Sub-partitioning
  within an original cell alone does not eliminate edges
  crossing other original cells; the global-component view is
  what licenses M4.  In practice strategy~(b) coincides with
  strategy~(a) on the global-component partition; the
  difference is presentational (sub-partitioning makes the
  component structure explicit before priority assignment).
\item[(c)] \emph{Boundary residual (literal-partition
  variant; default).}  Define a literal
  set-partition of $P$ into five buckets $\{R, C, D, S, B\}$
  by stipulation: classify each capability into the
  highest-priority cell among its membership-determined seeds
  (as in~(a) without the priority-extension absorption step),
  and place every \emph{boundary-crossing cooperative node}
  (cooperatives in $P$ whose participants span two or more
  original cells) into~$B$ rather than absorbing it into a
  participant's cell.  Cooperative nodes carry explicit
  $\volL$ weight $w(c_{\mathrm{coop}})$ from M5; the
  literal-partition assignment makes
  $\volL(B) = \sum_{z \in \mathrm{cross-cell\ coops}} w(z)$
  by direct M5 weight summation.  Subsumption-crossing
  edges, which carry no standalone $\volL$ weight, do not
  contribute additional bucket mass; their structural effect
  on additivity is what the literal-partition framing
  sidesteps.  The five-cell identity
  \[
    \volL(P) = \volL(R) + \volL(C) + \volL(D) + \volL(S) +
    \volL(B)
  \]
  follows from \emph{literal weight summation over the
  partitioned capability census}, \emph{without invoking
  M4-style poset-independence} between cells.  Under
  strategy~(c) we adopt the atomic-weight accounting convention
  $\volL(X) := \sum_{c \in X} w(c)$ for each cell~$X$ as an
  operational definition of cell-level $\volL$ values, distinct
  from $\volL$'s general M1--M6 poset-measure semantics on
  poset-independent capability subsets.  Cells partition $P$
  set-disjointly, so $\volL(P) = \sum_X \volL(X)$ follows by
  direct rearrangement of the atomic sum.  This is a strictly
  weaker license than~(a)'s M4-licensed identity: strategy
  (c) does not claim cross-cell poset-independence and does
  not certify M6-superadditivity properties of the partition;
  the cell-level $\volL$ values it produces are accounting
  aggregates, sufficient for the $\delta$-decomposition but
  not interchangeable with the M1--M6 measure on
  poset-independent components.
  What it provides is the additive identity needed for the
  $\delta$-decomposition, plus the diagnostic interpretation
  of $\delta_{\mathrm{boundary}}$ as ``cross-cell cooperative
  mass.''  The identity becomes $1 - \betaL =
  \delta_{\mathrm{restricted}} + \delta_{\mathrm{covered}} +
  \delta_{\mathrm{dormant}} + \delta_{\mathrm{residual}} +
  \delta_{\mathrm{boundary}}$ with all five terms well-defined
  as $\volL$-fractions.  Strategy~(c) is the paper's default
  and the variant the worked example uses
  (Section~\ref{sec:worked_example}): it
  preserves seed-cell attribution (dormant atoms stay in $D$
  rather than being absorbed into $C$ as under (a)) so that
  $\delta_{\mathrm{covered}}$ remains diagnostically
  interpretable as proxy-failure on observed capabilities and
  $\delta_{\mathrm{dormant}}$, $\delta_{\mathrm{residual}}$
  retain their literal seed-classification meanings, at the
  cost of foregoing the stronger M4-additivity guarantee that
  (a) provides.
\end{enumerate}
The paper's stated results assume strategy~(c) by default: the
literal-partition rule preserves the diagnostic-interpretation
discipline that the worked example demonstrates, at the
cost of a weaker (M5-only) additivity license.  Strategies (a)
and (b) remain available for deployments that prefer the
stronger M4-additivity guarantee and accept the
$\delta_{\mathrm{covered}}$ absorption ambiguity in exchange;
the practical prevalence of boundary-crossing poset edges is
small when the capability partition's priority classes are
themselves well-structured (e.g., restrictions are
poset-boundary closed because restricting a cooperative implies
restricting its participants, and restricting a capability
implies restricting capabilities subsumed by it in cost-aligned
posets), so the choice between (a)/(b) and (c) is largely a
matter of which interpretive discipline the deployment
requires.
\end{remark}

\begin{proposition}[Computability of Gap Decomposition,
given precomputed governance inputs]
\label{prop:gap_computable}
The complexity bound below is for the
\emph{cell-classification step alone}, given precomputed
governance inputs as listed.  Producing those inputs
(admissibility labels via institutional attestation,
residual classification by domain-expert review,
governance-set restriction set $\Xi$ and sufficiency
thresholds) is itself a non-trivial pipeline whose cost is
not part of this complexity claim.
Given
(i)~the aggregate sacrifice data (events in $[T - H, T]$),
(ii)~the restriction set $\Xi$,
(iii)~the residual class $\ResS$
(Definition~\ref{def:residual_class}),
(iv)~the capability census (enumeration of $P$) together with
the \emph{poset structure} $\mathcal{M}$: the
cooperative-composition hypergraph (cooperative nodes and their
participants) and the subsumption relation graph (directed
edges $c \succeq c'$),
(v)~an event-admissibility labelling designating which observed
events satisfy the theorem-admissibility requirements.  Because
Definition~\ref{def:gap_decomp} operates under Part~B
(the aggregate lower bound theorem~\citep[Revealed~Sacrifice]{lasser2026paper8a}), the full Part~B
admissibility set is required: S0, S1, S2, S4(a), S3, S5, the
genuine-exchange conditions (i)--(iv) of
the revealed-sacrifice event definition~\citep[Revealed~Sacrifice]{lasser2026paper8a}, \emph{and} within-bundle
poset-independence per the event-local bundle decomposition~\citep[Revealed~Sacrifice]{lasser2026paper8a}; and
(vi)~a closure policy from
Remark~\ref{rem:cell_cooperative_closure}, with strategy~(c)
(literal-partition with boundary-residual bucket) as the
default and (a)/(b) available as alternatives,
the gap decomposition (five-term under~(c), four-term under
(a)/(b)) is computable in time
$O\bigl(|P| + E + M \cdot \alpha(M)\bigr)$, where $E$ is the
number of observed (capability, event) membership records,
$M = |\mathcal{M}|$ is the total size of the poset structure
(cooperative hyperedges plus subsumption edges), and
$\alpha(\cdot)$ is the inverse Ackermann function from the
amortized analysis of union-find.  Since $\alpha(M)$ is
$\leq 5$ for any $M$ realizable in physical computation, this
is near-linear $\widetilde{O}(|P| + E + M)$ for practical
purposes.  The $O(|P|)$ term is per-individual-capability
classification; the $O(E)$ term is one-pass aggregation over
bundle-incidence records; the $O(M \cdot \alpha(M))$ term is
the closure-policy pass, which traverses cooperative and
subsumption edges jointly via amortized union-find.
No structural counterfactual computation is required (a
direct-measurement approach would require $O(|P|^2)$
counterfactual queries).
\end{proposition}

\begin{proof}
\emph{Pass 1 (individual-capability classification,
$O(|P|)$).}
Maintain a hash table from capabilities to their current class
labels, initialized in $O(|P|)$ time.  Sweep the individual
capabilities in $P$ once: for each $c \in P$, check $c \in \Xi$
(restricted table lookup, $O(1)$); if not, check $c \in \ResS$
(residual classification lookup, $O(1)$); if in neither, mark as
dormant.  This gives an initial classification in $O(|P|)$.

\emph{Pass 2 (event-sweep, $O(E)$).}
Iterate through each (capability, event) incidence record.
For each record $(c, n)$ with event $n$ admissibility-labelled
satisfying the \emph{full Part~B admissibility set} of input
(v)---S0, S1, S2, S4(a), S3, S5, the genuine-exchange
conditions, and within-bundle poset-independence of $Y_n$ per
the event-local bundle decomposition~\citep[Revealed~Sacrifice]{lasser2026paper8a}---accumulate the
event's $\volRlower$ contribution to $c$ regardless of $c$'s
priority class (so that $\volRlower(R)$ for restricted
capabilities and $\volRlower(C)$ for covered capabilities are
both maintained correctly per
Definition~\ref{def:gap_decomp}'s $\delta$ formulas, which
subtract $\volRlower(R)$ and $\volRlower(C)$ respectively).
The class label set in Pass 1 is preserved
(restricted $\to$ restricted, residual $\to$ residual, dormant
$\to$ covered when receiving evidence); only the
$\volRlower$ accumulation runs unconditionally.  Events failing
any component of the Part~B admissibility set do not contribute
coverage; in particular, events with internally cooperative
bundles are filtered out by the within-bundle
poset-independence check.  This is $O(E)$ total work, one
constant-time update per record.

\emph{Pass 3 (closure-policy pass, $O(M)$).}
Under default strategy~(c), walk the cooperative hyperedges
identifying boundary-crossing cooperative nodes (cooperatives
whose participants span two or more original cells) and place
each into the boundary-residual bucket~$B$; non-boundary-crossing
cooperatives remain in their participants' cell.  Subsumption
edges contribute no $\volL$ weight and are therefore not
processed; the literal-partition framing sidesteps their
additivity contribution.  Under alternative strategy~(a), walk
the combined
cooperative/subsumption graph via union-find, identifying
poset-connected components: two nodes are in the same component
if connected by any chain of cooperative-participation or
subsumption edges.  For each component, determine its highest-
priority cell (restricted $>$ covered $>$ dormant $>$ residual)
among member cells, and place the entire component in that
cell.  The cells become both cooperative-closed and
subsumption-closed (hence poset-boundary closed per
Definition~\ref{def:gap_decomp}) by construction.  Union-find
gives amortized $O(M \cdot \alpha(M))$ complexity, essentially
linear in $M$.  Under strategy~(b), the pass is replaced
by the sub-partition construction over the same combined graph,
with the same complexity modulo constants.  Strategy~(c)'s
single hyperedge sweep is $O(M)$ outright (no union-find
needed since no transitive-closure of cell membership is
performed).

\emph{Aggregation.}
Each class's aggregate $\volL$ and $\volRlower$ is maintained
incrementally across the three passes; computing the four
quotients from the running totals at the end of Pass~3 is
$O(|\mathrm{cells}|) = O(1)$ since there is a constant number
of cells, but the running totals themselves are updated as
nodes move between cells under closure (Pass~3), so the
total accounting cost is folded into Passes~1--3 rather than
literally appended.  The counterfactual
comparison---removing capability $c$ and recomputing global
reachability---is not performed; only membership checks,
aggregation, and closure-policy lookups occur.  A
direct-measurement approach requiring $O(|P|^2)$ counterfactuals
is avoided because covered/dormant classification here uses
bundle-appearance evidence, not marginal-contribution evidence.
\end{proof}

\begin{remark}[Admissibility-labelling is an external input]
\label{rem:admissibility_external}
S0 (calibration) and the interior component of S1 are not
mechanically verifiable from the observable event fields alone.
Input~(v) of Proposition~\ref{prop:gap_computable} represents an
external admissibility layer: events must be labelled admissible
(by trust-weighted attestation, auditor sampling, or explicit
agent commitment) before the decomposition can classify
capabilities into ``covered.''  Under committed-mode operation
(\citep[Revealed~Sacrifice]{lasser2026paper8a}), the labelling is carried by the
commitment protocol's proof attachment; in unrestricted mode, the
framework must rely on governance processes to supply it.  The
$O(|P|)$ complexity is the \emph{classification} cost given the
labelling, not the cost of producing the labelling itself.
\end{remark}

\begin{remark}[Part~A versus Part~B granularity]
Under Part~A alone (category-level aggregates, requiring S0,
S1, S2, S4(a), and the genuine-exchange conditions of the
revealed-sacrifice event definition in the companion
paper~\citep[Revealed~Sacrifice]{lasser2026paper8a}), the
decomposition operates at the category level: each partition
component $C_j$
(the category partition~\citep[Revealed~Sacrifice]{lasser2026paper8a}) is classified as
restricted, covered, dormant, or residual.  Under committed data
(\citep[Revealed~Sacrifice]{lasser2026paper8a}), S2 splits into
two components with different attestation routes: (i)~the
\emph{objective} side---that $\Delta\volL(X)$ is correctly
computed from the surrendered benchmarked capability via the
public benchmark function---is verified by the commitment
layer's ZK proof on the magnitude commitment; (ii)~the
\emph{subjective} side---that the agent's valuation $\Ui(X)$
of the surrendered capability matches $\Delta\volL(X)$
in the calibration scale---is not ZK-verifiable and must be
supplied by external attestation
(see the external-attestation remark in the companion
paper~\citep[Revealed~Sacrifice]{lasser2026paper8a}).
The committed-mode decomposition therefore certifies the
objective S2 component via the commitment layer and assumes the
subjective S2 component via the attestation layer.  The
per-capability decomposition is
the refinement available when Part~B assumptions hold.
\end{remark}

\subsection{Need-sufficiency diagnostic (companion metric)}
\label{sec:need_diagnostic}

The gap decomposition above tracks want-pursuit: how much of the
collective's above-sufficiency capability-space is evidenced by
voluntary sacrifice.  The need-sufficiency diagnostic is the
parallel metric for below-sufficiency: how efficiently does the
collective satisfy needs, and is the satisfaction downstream-safe?

\begin{definition}[Need-Access Cost]
\label{def:nac}
The need-access cost aggregate for trade window $[T - H, T]$ is:
\[
  \NAC(T) \;=\; \sum_{n \in F_{S1}}
  \Delta\volL(X_n),
\]
where $F_{S1}$ is the set of events failing S1 due to
below-sufficiency (Definition~\ref{def:s1_sufficiency}).
$\NAC$ measures the total benchmarked capability surrendered for
need-satisfaction---purchasing power and time spent reaching
baseline functionality rather than pursuing wants.

$\NAC$ does \emph{not} enter
the aggregate lower bound of the companion
paper~\citep[Revealed~Sacrifice]{lasser2026paper8a} or
$\betaL$.  It is a separate diagnostic that tracks the cost side
of the agent's capability budget.
\end{definition}

\begin{remark}[$\NAC$ is an upper bound on need-access cost]
\label{rem:nac_upper_bound}
Because the whole-event filter sends the entire
$\Delta\volL(X_n)$ of mixed need/want events into $\NAC$,
$\NAC$ overestimates the true need-access cost by including any
attached want-pursuit spend.  $\NAC$ should be read as an
\emph{upper bound} on need-access cost, not a precise
measurement.

\emph{Looseness magnitude depends on the deployment's mixed-event
prevalence.}  In economies where most below-sufficiency events
involve pure-need bundles (a clinic visit purely for medical care,
a bottle of water purely for hydration), $\NAC$ is close to the
true need-access cost.  In economies where below-sufficiency
events typically carry attached want-pursuit (a grocery basket
that contains both staple food at below-sufficiency-cost wages and
discretionary items, a transit pass that provides both
need-satisfaction and recreational mobility), $\NAC$ can
substantially overstate the true need-access cost---potentially
by factors of two or more in dense mixed-bundle markets.  The
diagnostic interpretation of ``high $\NAC$ alarms'' in
Proposition~\ref{prop:combined_diagnostic}'s quadrants is
therefore sensitive to the deployment's mixed-bundle prevalence:
a low-mixed-prevalence deployment can read $\NAC$ thresholds
nearly literally; a high-mixed-prevalence deployment should
treat $\NAC$ thresholds as conservative upper-bound triggers
that may fire on inflated values rather than on true
need-access cost increases.  A bundle-decomposition
extension that splits $\Delta\volL(X_n)$ into need-satisfying
and want-pursuing components via the event-local bundle
decomposition of the companion
paper~\citep[Revealed~Sacrifice]{lasser2026paper8a} would tighten $\NAC$ to
the need-component aggregate and is the natural follow-up; it
is not part of this paper's $\NAC$ definition because the
event-local decomposition coefficients $\alpha_{n,c}$ are not
generally available at the granularity required without
additional governance attestation.
\end{remark}

\begin{definition}[Need-Satisfaction Efficiency]
\label{def:nse}
For each need $N$ identified in the poset
(Definition~\ref{def:need_bundle}) and each agent~$i$, the
per-agent need-satisfaction efficiency is a triple:
\[
  \NSE(i, N) \;=\; (\mathrm{sufficiency\_status},\;
  \mathrm{downstream\_safety},\; \mathrm{access\_cost}),
\]
where:
\begin{itemize}
\item $\mathrm{sufficiency\_status} \in
  \{\text{below}, \text{met}, \text{exceeded}\}$;
\item $\mathrm{downstream\_safety} \in
  \{\text{safe}, \text{degrading}, \text{unknown}\}$;
\item $\mathrm{access\_cost} = \sum \Delta\volL(X_n)$ for
  below-sufficiency events by agent~$i$ targeting need~$N$.
\end{itemize}
The collective-level aggregate $\NSE(T)$ reports the distribution
of (agent, need) pairs across sufficiency statuses and flags
degrading pathways.
\end{definition}

\begin{remark}[Observability under commitment]
When sacrifice events are committed
(\citep[Revealed~Sacrifice]{lasser2026paper8a}), per-agent $\NSE(i, N)$ is not
recoverable.  The committed-mode fallback is a population-level
proxy: $\NSE_{\mathrm{proxy}}(T)$ reports $\NAC(T)$ by need
category and counts of below-sufficiency events.  Full $\NSE(T)$
requires agents to voluntarily attest pathways or the framework to
operate in non-committed mode.
\end{remark}

\begin{proposition}[Combined Diagnostic]
\label{prop:combined_diagnostic}
$\volRlower$ and $\NAC$ are not directly commensurable as
written: $\volRlower$ is a poset-measure quantity bounded above
by $\volL(P)$, while $\NAC$ is a cumulative event aggregate of
$\Delta\volL(X_n)$ summands that can grow without bound as more
below-sufficiency events accumulate within the window.  For
diagnostic comparison, we use the normalised pair
$(\volRlower / \volL(P),\, \NAC / B_W)$ where $B_W$ is a
domain-supplied total benchmarked-budget scale for window~$W$
(e.g., aggregate purchasing power exchanged in $W$, or a
normalised-per-agent equivalent).  ``High'' and ``low'' below
refer to these normalised quantities crossing
governance-supplied thresholds, not to the raw aggregates.  The
framework reports the triple
$(\volRlower, \NAC, \NSE)$ as a joint diagnostic:
\begin{enumerate}
\item[(a)] \textbf{High $\volRlower$ + low $\NAC$:} Healthy.
  The collective exercises rich capabilities above sufficiency
  and needs are cheap to meet.
\item[(b)] \textbf{Low $\volRlower$ + high $\NAC$:} Proxy
  failure candidate.  $\volL$ may be high, but agents are
  spending most of their benchmarked capability budget on
  basic need-satisfaction.  The sacrifice channel excludes the
  H1+H2 below-sufficiency sub-class characterized by
  Definition~\ref{def:s1_sufficiency} (retention degrading and
  no third non-degrading alternative, i.e., the sub-class where
  S1 unambiguously fails), so a low $\volRlower$
  alongside a high $\NAC$ is \emph{consistent with} low
  above-sufficiency exercise; the gap-decomposition's
  alarm-interpretation caveats
  (Remark~\ref{rem:alarm_interpretation}) still apply (the low
  $\volRlower$ may also reflect bound-looseness, sparse
  observation, or admissibility filtering).  $\NAC$ surfaces a
  candidate explanation for the low $\volRlower$.
\item[(c)] \textbf{Low $\volRlower$ + low $\NAC$:} Dormancy,
  residual, or theorem-inadmissible activity.  The gap
  decomposition (Definition~\ref{def:gap_decomp}) distinguishes
  dormant from residual.
\item[(d)] \textbf{High $\volRlower$ + high $\NAC$:} Mixed
  regime.  $\NSE$ distinguishes ``expensive but safe'' from
  ``expensive and degrading.''
\end{enumerate}
\end{proposition}

\begin{remark}[Parallel structure, not nesting]
The need-sufficiency diagnostic is parallel to the gap
decomposition, not nested within it, but operates on different
formal objects.  The gap decomposition is a \emph{capability
partition}; $\NAC$ is an \emph{event aggregate}; $\NSE$ is a
\emph{per-need state assessment}.  The two diagnostics cover
complementary portions of the sacrifice event stream:
S1-satisfying events (that also satisfy S0 and S2) enter the gap
decomposition via the aggregate lower bound theorem~\citep[Revealed~Sacrifice]{lasser2026paper8a}; events
failing S1 due to below-sufficiency enter $\NAC$.  Events failing
S1 for non-need reasons (coercion, duress) are excluded from
both diagnostics and constitute an undiagnosed event-stream
residual.
\end{remark}
